\section{El protocolo de comunicación sensor--plataforma}
\label{sec:protocolo}

Esta sección desarrolla la decisión técnica de mayor peso del proyecto y responde a la instrucción
explícita del profesor guía: determinar cómo funciona el protocolo, dejarlo definido y justificar por
qué se emplea este y no otro.

\subsection{Qué es Contact ID}

\subsubsection{Origen y posición en la familia de estándares SIA}

Contact ID es un protocolo de reporte de eventos creado por Ademco para que los paneles de alarma
comuniquen sucesos a las centrales de monitoreo. Fue formalizado por la \emph{Security Industry
Association} como estándar \textbf{DC-05}, cuya edición vigente es DC-05-2016.

Conviene distinguir tres estándares SIA que suelen confundirse, porque \textbf{no compiten entre sí sino
que se apilan}.

\begin{longtable}{@{}L{3.4cm}L{2.6cm}L{7.7cm}@{}}
\filaencab \encab{Estándar} & \encab{Capa} & \encab{Qué define} \\ \midrule \endhead
\textbf{DC-05} (Contact ID) & \textbf{Aplicación} & El formato del mensaje de evento: qué campos posee y qué significa cada código \\
DC-03 (SIA) & Aplicación & Formato alternativo, más expresivo pero más extenso \\
\textbf{DC-09} & \textbf{Transporte} & Cómo se transportan mensajes DC-05 o DC-03 sobre IP (TCP o UDP), con cifrado AES opcional \\
\bottomrule
\caption{Estándares SIA relevantes y su posición en la pila de protocolos.}
\label{tab:sia}
\end{longtable}

\begin{cajaclave}
La elección no se plantea entre Contact ID y DC-09. Contact ID es el formato del mensaje y DC-09 su
transporte moderno sobre redes IP: \textbf{se emplean conjuntamente}.
\end{cajaclave}

La revisión \textbf{DC-09-2026}, publicada en marzo de 2026 con aprobación ANSI, incorpora dos elementos
directamente pertinentes a este trabajo: auto-comisionamiento de dispositivos y \textbf{rotación de
claves de cifrado}. Ambos corresponden a preocupaciones de administración de red ---aprovisionamiento
seguro y gestión del ciclo de vida de credenciales--- que constituyen el eje del proyecto.

\subsubsection{Gramática del mensaje}

Un mensaje Contact ID es una secuencia de \textbf{16 dígitos}.

\begin{figure}[htbp]
\centering
\begin{tikzpicture}

% ---- campos del mensaje ----
\node[campo, minimum width=17mm, fill=azuloscuro!8] (acct) {1234};
\node[campo, minimum width=11mm, right=0pt of acct, fill=grisclaro] (mt) {18};
\node[campo, minimum width=8mm,  right=0pt of mt,   fill=ambar!15]  (q)   {1};
\node[campo, minimum width=14mm, right=0pt of q,    fill=verdeacento!12] (xyz) {381};
\node[campo, minimum width=11mm, right=0pt of xyz,  fill=grisclaro] (gg)  {00};
\node[campo, minimum width=14mm, right=0pt of gg,   fill=verdeacento!12] (ccc) {007};
\node[campo, minimum width=8mm,  right=0pt of ccc,  fill=grisclaro] (s)   {B};

% ---- nombres de campo encima ----
\foreach \n/\l in {acct/ACCT, mt/MT, q/Q, xyz/XYZ, gg/GG, ccc/CCC, s/S}
  \node[font=\scriptsize\bfseries\color{azuloscuro}, above=1pt of \n] {\l};

% ---- descripciones debajo, escalonadas ----
\node[nota, below=17mm of s.south, xshift=6mm, anchor=west]   (ds)   {checksum};
\node[nota, below=24mm of ccc.south, xshift=6mm, anchor=west] (dccc) {zona 007: \textbf{el dispositivo}};
\node[nota, below=31mm of gg.south, xshift=6mm, anchor=west]  (dgg)  {grupo o partición};
\node[nota, below=38mm of xyz.south, xshift=6mm, anchor=west] (dxyz) {código de evento 381: \textbf{pérdida de supervisión RF}};
\node[nota, below=45mm of q.south, xshift=6mm, anchor=west]    (dq)   {calificador 1: \textbf{nuevo evento} (3 = restauración)};
\node[nota, below=52mm of mt.south, xshift=6mm, anchor=west]  (dmt)  {tipo de mensaje, siempre 18};
\node[nota, below=59mm of acct.south, xshift=6mm, anchor=west](dacct){cuenta 1234: \textbf{el sitio}};

\foreach \c/\d in {s/ds, ccc/dccc, gg/dgg, xyz/dxyz, q/dq, mt/dmt, acct/dacct} {
  \draw[flecha, dashed, line width=0.4pt] (\c.south) |- (\d.west);
}

% ---- lectura ----
\node[cajaacento, below=6mm of dacct.south west, anchor=north west, text width=12.5cm,
      fill=verdeacento!6, font=\footnotesize\itshape]
  {«En la cuenta 1234 se produjo un nuevo evento de pérdida de supervisión RF
    en el dispositivo 007.»};

\end{tikzpicture}
\caption{Estructura del mensaje Contact ID: 16 dígitos distribuidos en siete campos.}
\label{fig:mensaje}
\end{figure}


\begin{longtable}{@{}L{1.8cm}C{1.6cm}L{10.3cm}@{}}
\filaencab \encab{Campo} & \encab{Dígitos} & \encab{Contenido} \\ \midrule \endhead
\textbf{ACCT} & 4 & Número de abonado o cuenta \\
\textbf{MT} & 2 & Tipo de mensaje. Siempre 18 (preferido) o 98 (opcional) \\
\textbf{Q} & 1 & Calificador: \textbf{1} nuevo evento o apertura; \textbf{3} restauración o cierre; \textbf{6} estado previo aún presente \\
\textbf{XYZ} & 3 & Código de evento \\
\textbf{GG} & 2 & Grupo o partición. El valor 00 indica que no aplica \\
\textbf{CCC} & 3 & Número de zona, para eventos, o de usuario, para aperturas y cierres \\
\textbf{S} & 1 & Checksum \\
\bottomrule
\caption{Campos del mensaje Contact ID.}
\label{tab:campos}
\end{longtable}

\begin{cajaclave}[title={El calificador convierte el protocolo en una máquina de estados}]
El campo \texttt{Q} es lo que transforma un aviso simple en una máquina de estados: el valor 1 abre una
condición y el valor 3 la cierra. Aplicado a la gestión de sensores, el protocolo distingue nativamente
entre \enquote{el nodo 007 dejó de reportar} (\texttt{Q=1}, código 381) y \enquote{el nodo 007 volvió a
reportar} (\texttt{Q=3}, código 381). \textbf{La plataforma no necesita inferir el término de una
condición: el protocolo lo declara.}
\end{cajaclave}

\subsubsection{Alfabeto y cálculo del checksum}

Contact ID no emplea el alfabeto decimal habitual, porque su transporte original fue \textbf{DTMF} ---los
tonos del teclado telefónico---, donde no existe un tono para el cero. El alfabeto consta de 15 símbolos.

\begin{center}
\begin{tabular}{@{}l|ccccccccc|c|ccccc@{}}
\textbf{Símbolo} & 1 & 2 & 3 & 4 & 5 & 6 & 7 & 8 & 9 & \textbf{0} & \textbf{B} & \textbf{C} & \textbf{D} & \textbf{E} & \textbf{F} \\
\midrule
\textbf{Valor} & 1 & 2 & 3 & 4 & 5 & 6 & 7 & 8 & 9 & \textbf{10} & \textbf{11} & \textbf{12} & \textbf{13} & \textbf{14} & \textbf{15} \\
\end{tabular}
\captionof{table}{Alfabeto de Contact ID y valor numérico de cada símbolo.}
\label{tab:alfabeto}
\end{center}

El checksum $S$ es el símbolo que satisface:

\begin{equation}
\left(\sum_{i=1}^{15} v(d_i) + v(S)\right) \bmod 15 = 0
\end{equation}

donde $v(d)$ denota el valor del dígito $d$ según la tabla~\ref{tab:alfabeto}.

\paragraph{Ejemplo trabajado.} Cuenta 1234, nuevo evento, código 381 (pérdida de supervisión RF), grupo
00, zona 007.

\begin{center}
\begin{tabular}{@{}l|ccccccccccccccc|c@{}}
\textbf{Dígito} & 1 & 2 & 3 & 4 & 1 & 8 & 1 & 3 & 8 & 1 & 0 & 0 & 0 & 0 & 7 & $\Sigma$ \\
\midrule
\textbf{Valor} & 1 & 2 & 3 & 4 & 1 & 8 & 1 & 3 & 8 & 1 & \textbf{10} & \textbf{10} & \textbf{10} & \textbf{10} & 7 & \textbf{79} \\
\end{tabular}
\end{center}

Dado que $79 \bmod 15 = 4$, falta $15 - 4 = 11$ para completar el múltiplo. El símbolo de valor 11 es
\textbf{B}. El mensaje completo es \texttt{123418138100007B}, verificable: $79 + 11 = 90$ y
$90 \bmod 15 = 0$.

\begin{cajaclave}[title={Verificación empírica}]
El cálculo no se asume. La herramienta \texttt{plataforma/herramientas/contactid.py}, desarrollada como
parte de este análisis, genera y valida mensajes, y comprueba que la alteración de un solo dígito
invalida el checksum. Todos los ejemplos y tamaños de esta sección fueron calculados con ella.
\end{cajaclave}

\textbf{Capacidad de detección de errores.} El checksum módulo 15 detecta cualquier error de un solo
dígito, pero \textbf{no detecta transposiciones}: intercambiar dos dígitos no altera la suma. Es una
limitación que debe documentarse. Sobre un transporte con integridad propia ---LoRaWAN aporta un código
de integridad de 4 bytes, y TLS su propio mecanismo--- el checksum opera como verificación de sanidad de
la capa de aplicación, no como único mecanismo de integridad.

\subsection{El proceso de transmisión}

\subsubsection{Transporte histórico: DTMF sobre línea telefónica}

\begin{figure}[htbp]
\centering
\begin{tikzpicture}[font=\footnotesize]

\def\lx{0}
\def\rx{9.2}

\node[cajaacento, minimum width=26mm, font=\footnotesize] (panel) at (\lx,0) {\textbf{Panel de alarma}};
\node[cajaacento, minimum width=30mm, font=\footnotesize] (recep) at (\rx,0) {\textbf{Central receptora}};

\draw[line width=0.6pt, color=grislinea, dashed] (panel.south) -- (\lx,-7.0);
\draw[line width=0.6pt, color=grislinea, dashed] (recep.south) -- (\rx,-7.0);

\draw[flecha] (\lx,-1.0) -- node[nota, above, font=\scriptsize] {descuelga y marca} (\rx,-1.0);

\draw[flecha] (\rx,-2.2) -- node[nota, above, font=\scriptsize] {\textit{handshake}: tonos 1400 Hz + 2300 Hz} (\lx,-2.2);
\node[nota, font=\scriptsize\itshape, below=0.5mm of {(4.6,-2.2)}] {«hablo Contact ID»};

\draw[flecha, line width=1.2pt, color=verdeacento] (\lx,-3.7) --
  node[nota, above, font=\scriptsize\bfseries\color{verdeacento}] {16 dígitos en DTMF ($\approx$ 1,6 s)} (\rx,-3.7);

\draw[flecha, color=ambar, line width=1pt] (\rx,-4.9) --
  node[nota, above, font=\scriptsize\bfseries\color{ambar}] {\textit{kissoff}: tono 1400 Hz} (\lx,-4.9);
\node[nota, font=\scriptsize\itshape, below=0.5mm of {(4.6,-4.9)}] {«recibido, checksum correcto»};

\draw[flecha] (\lx,-6.2) -- node[nota, above, font=\scriptsize] {cuelga} (\rx,-6.2);

\node[cajaambar, below=4mm of {(4.6,-7.0)}, text width=11.6cm, fill=ambar!6,
      font=\scriptsize]
  {El \textbf{\textit{kissoff}} solo se emite si el checksum es correcto. Si el panel no lo recibe,
   \textbf{reintenta}: el protocolo incorpora confirmación explícita a nivel de aplicación,
   independiente del transporte.};

\end{tikzpicture}
\caption{Secuencia de transmisión original de Contact ID sobre línea telefónica.}
\label{fig:dtmf}
\end{figure}


El elemento relevante es el \textbf{kissoff}: la receptora lo emite únicamente si el checksum resulta
correcto; en su ausencia, el panel reintenta. El protocolo incorpora por tanto \textbf{confirmación
explícita a nivel de aplicación}, independiente del transporte. Ese comportamiento se conserva como
requisito de diseño en el perfil propuesto.

\subsubsection{El caso de las alarmas monitoreadas}

El profesor guía señaló como referencia los sensores de empresas de alarmas monitoreadas y su
comunicación mediante llamada celular. Conviene precisar el mecanismo, porque explica el diseño del
formato.

\begin{cajaclave}[title={La comunicación celular es un cambio de transporte, no de protocolo}]
Los comunicadores celulares del sector convierten la marcación DTMF de línea fija en marcación DTMF
sobre red móvil, \textbf{conservando el mensaje Contact ID intacto}. El protocolo de aplicación no ha
cambiado: cambió el medio.

Esto explica por qué el mensaje consta de 16 dígitos y no de un objeto estructurado: debía caber en
tonos telefónicos sobre una línea analógica. \textbf{Y esa restricción de hace cuatro décadas es
precisamente la que hoy lo hace apropiado para LoRa}: el protocolo fue diseñado para un canal estrecho,
poco fiable y costoso por transmisión, que es la misma descripción del enlace de un nodo LoRa
alimentado a batería.
\end{cajaclave}

\begin{figure}[htbp]
\centering
\resizebox{\linewidth}{!}{%
\begin{tikzpicture}[font=\scriptsize]

\node[caja, text width=2.3cm, align=center, draw=ambar, line width=1pt] (sens) at (0,0)
  {\textbf{Sensores}\\ inalámbricos\\ a batería};
\node[caja, text width=2.3cm, align=center, draw=ambar, line width=1pt] (panel) at (4.6,0)
  {\textbf{Panel} o\\ comunicador};
\draw[flecha] (sens.east) -- node[nota, below=2pt, font=\scriptsize, align=center] {radio\\propietaria} (panel.west);

% los tres transportes, como cajas apiladas (ningún texto sobre líneas)
\node[caja, text width=4.9cm, align=left, draw=grislinea] (t1) at (9.0, 1.7)
  {\textbf{1. Línea telefónica fija}\\ tonos DTMF \ \textit{(histórico)}};
\node[caja, text width=4.9cm, align=left, draw=grislinea] (t2) at (9.0, 0)
  {\textbf{2. Red celular}\\ tonos DTMF sobre 2G/3G \ \textit{(en retirada)}};
\node[caja, text width=4.9cm, align=left, draw=verdeacento, line width=1.1pt, fill=verdeacento!6] (t3) at (9.0,-1.7)
  {\textbf{3. Internet}\\ SIA DC-09 con AES \ \textit{(el adoptado)}};

\node[caja, text width=2.4cm, align=center, draw=azuloscuro, line width=1.1pt, fill=azuloscuro!5]
  (arc) at (14.4,0) {\textbf{Central}\\ \textbf{receptora}\\ (ARC / CMS)};

\foreach \t in {t1,t2,t3}{
  \draw[flecha] (panel.east) -- (\t.west);
  \draw[flecha] (\t.east) -- (arc.west);
}

\node[notaclave, text width=14.2cm, align=left, font=\footnotesize] at (7.2,-4.0)
  {\textbf{En los tres casos el mensaje de aplicación es el mismo: Contact ID.} La «comunicación celular»
   es un cambio de transporte, no de protocolo. El mensaje tiene 16 dígitos porque debía caber en tonos
   telefónicos sobre una línea analógica, \textbf{y esa restricción de hace cuatro décadas es la que hoy
   lo hace apropiado para LoRa}, un canal igualmente estrecho.};
\end{tikzpicture}}
\caption{Cadena de comunicación de una alarma monitoreada: el formato del mensaje se conserva a través
de tres generaciones de transporte.}
\label{fig:alarmas}
\end{figure}


Un caso ilustrativo: Verisure Securitas Direct opera con red celular como canal principal e incorporó
Sigfox como canal secundario en su parque instalado. Cuando el panel detecta pérdida de señal celular
---indicio de un intento de interferencia--- emite la alerta por la red alternativa para garantizar
continuidad del servicio. Dos observaciones se desprenden:

\begin{enumerate}
\item Es multiprotocolo por diseño y por la razón correcta: un operador que compromete continuidad de
servicio no se ata a un solo transporte. Refuerza el requisito de no cerrarse a una única tecnología de
radio.
\item La condición detectada ---pérdida del enlace por interferencia--- corresponde al código 344 de
Contact ID. El vocabulario del protocolo ya contempla ese evento.
\end{enumerate}

\begin{cajaclave}[title={El planteamiento del trabajo, en una formulación}]
La industria de la seguridad electrónica resolvió hace décadas el problema de supervisar flotas de
sensores inalámbricos alimentados a batería, y estandarizó su vocabulario en un protocolo
\textbf{abierto} (ANSI/SIA DC-05). Sin embargo, lo aplica dentro de \textbf{sistemas cerrados y
verticalmente integrados}: los sensores solo se comunican con el panel del proveedor, el panel solo
reporta a su central, y el cliente no puede retirar sus datos ni incorporar equipamiento de terceros.
Esto constituye precisamente las causas C2 y C3 del problema definido.

\textbf{Este trabajo toma ese protocolo estándar y lo emplea como contrato de interoperabilidad de una
plataforma abierta}, extendiendo su transporte a LoRaWAN, WiFi/HTTP y MQTT.
\end{cajaclave}

\subsubsection{Transporte adoptado: SIA DC-09 sobre IP}

DC-09 encapsula el mensaje DC-05 en una trama sobre TCP o UDP, identificando el formato mediante un
token: \texttt{ADM-CID} para Contact ID sin cifrar y \texttt{*ADM-CID} para su versión cifrada. Admite
\textbf{AES-128, AES-192 y AES-256} con clave precompartida.

Esto resuelve la exigencia de seguridad de forma nativa: no es necesario diseñar una capa de cifrado
propia, y el algoritmo resulta coherente con el que LoRaWAN emplea de fábrica.

El transporte DTMF sobre llamada celular se descarta porque las redes de segunda y tercera generación
están siendo desactivadas, y la transmisión de tonos depende de un canal conmutado que ya no está
garantizado.

\subsection{Vocabulario de eventos aplicable a la gestión de sensores}

\begin{longtable}{@{}C{1.6cm}L{5.4cm}L{6.7cm}@{}}
\filaencab \encab{Código} & \encab{Significado} & \encab{Uso en la plataforma} \\ \midrule \endhead
\textbf{381} & Pérdida de supervisión RF & El nodo dejó de reportar; con \texttt{Q=3}, se restableció \\
\textbf{384} & Batería baja en zona RF & Mantenimiento preventivo del nodo \\
\textbf{383} & Manipulación de sensor & Intervención física del dispositivo \\
\textbf{380} & Problema o avería de sensor & Falla del elemento sensor \\
\textbf{344} & Interferencia en receptor RF & Degradación del medio radioeléctrico \\
\textbf{302} & Batería baja del sistema & Alimentación del gateway \\
\textbf{301} & Pérdida de alimentación de red & Corte eléctrico en el sitio \\
\textbf{350} & Problema de comunicación & Falla del enlace del gateway \\
\textbf{602} & Reporte de prueba periódico & Confirma que el sistema opera \\
\textbf{603} & Transmisión RF periódica & Confirmación periódica de nodo por radio \\
\midrule
\textbf{158} & Temperatura alta & Cámara de frío, invernadero, sala de equipos \\
\textbf{159} & Temperatura baja & Riesgo de helada, cadena de frío \\
\textbf{154} & Fuga de agua & Estanques, riego, sala técnica \\
\textbf{161} & Pérdida de flujo de aire & Ventilación y climatización \\
\textbf{134} & Intrusión entrada/salida & Apertura de puerta de bodega \\
\textbf{132} & Intrusión perimetral & Perímetro de predio \\
\textbf{400} & Apertura o cierre & Acceso autorizado, con usuario en \texttt{CCC} \\
\textbf{570} & Zona anulada & Nodo puesto en mantenimiento \\
\bottomrule
\caption{Códigos de evento DC-05 aplicables a la gestión de sensores.}
\label{tab:codigos}
\end{longtable}

\begin{cajaclave}[title={Observación central del análisis}]
Los códigos 381, 384, 344, 602 y 603 constituyen, en conjunto, un \textbf{vocabulario normalizado de
salud de enlace de dispositivos inalámbricos} que existe desde antes de que se acuñara el término
\enquote{Internet de las Cosas}. La industria de seguridad electrónica resolvió el problema de
supervisar millares de nodos de radio alimentados a batería décadas atrás y estandarizó su vocabulario.
Este trabajo lo reutiliza en lugar de proponer uno nuevo.
\end{cajaclave}

\subsection{Evaluación de alternativas}

\subsubsection{Criterios y ponderación}

\begin{longtable}{@{}C{1.4cm}L{4.6cm}C{1.4cm}L{5.9cm}@{}}
\filaencab \encab{} & \encab{Criterio} & \encab{Peso} & \encab{Fundamento del peso} \\ \midrule \endhead
\textbf{C-A} & Compacidad del mensaje & 25\% & Restricción física: el payload útil de LoRaWAN puede descender a 11 bytes. Un protocolo que no quepa queda descartado por física, no por preferencia \\
\textbf{C-B} & Normalización semántica de eventos & 20\% & Es la causa C1. Un protocolo sin vocabulario cerrado reproduce el problema en lugar de resolverlo \\
\textbf{C-C} & Interoperabilidad con infraestructura instalada & 15\% & El problema se define sobre organizaciones que ya poseen equipamiento \\
\textbf{C-D} & Seguridad nativa del transporte & 15\% & Requisito de entorno de producción y eje del enfoque de redes \\
\textbf{C-E} & Soporte de telemetría con valor y unidad & 10\% & Necesario, pero resoluble mediante un plano complementario \\
\textbf{C-F} & Costo de implementación en nodo restringido & 15\% & Los nodos son microcontroladores con memoria de kilobytes y presupuesto energético estricto \\
\bottomrule
\caption{Criterios de evaluación del protocolo y su ponderación.}
\label{tab:criteriosprotocolo}
\end{longtable}

\subsubsection{Matriz comparativa}

\begin{longtable}{@{}L{4.6cm}C{1.2cm}C{1.2cm}C{1.2cm}C{1.2cm}C{1.2cm}C{1.2cm}C{1.5cm}@{}}
\filaencab \encab{Alternativa} & \encab{C-A} & \encab{C-B} & \encab{C-C} & \encab{C-D} & \encab{C-E} & \encab{C-F} & \encab{Total} \\
\midrule \endhead
MQTT + JSON ad hoc & 1 & 1 & 3 & 4 & 5 & 3 & 2,45 \\
CoAP (RFC 7252) & 4 & 1 & 2 & 4 & 4 & 3 & 2,95 \\
SenML (RFC 8428) & 3 & 3 & 2 & 2 & 5 & 3 & 2,90 \\
Sparkplug B & 2 & 4 & 3 & 4 & 5 & 2 & 3,15 \\
LwM2M (OMA) & 3 & 4 & 3 & 5 & 5 & 2 & 3,55 \\
CayenneLPP & 5 & 2 & 2 & 3 & 4 & 5 & 3,55 \\
SIA DC-03 & 3 & 5 & 5 & 4 & 4 & 3 & 3,95 \\
\textbf{Contact ID + DC-09} & \textbf{5} & \textbf{5} & \textbf{5} & \textbf{4} & \textbf{4} & \textbf{5} & \textbf{4,75} \\
\bottomrule
\caption{Matriz comparativa ponderada de alternativas de protocolo. Escala de 1 a 5. La puntuación de Contact ID en el criterio C-E se corrige de 1 a 4 conforme a la sección~\ref{sec:extdata}.}
\label{tab:matrizprotocolo}
\end{longtable}

\subsubsection{Justificación de las puntuaciones determinantes}

\textbf{MQTT con JSON ad hoc (2,45): constituye el problema, no la solución.} Obtiene 1 en normalización
semántica porque cada implementación define sus propios nombres de campo, que es exactamente la causa
C1. Obtiene 1 en compacidad: un objeto JSON típico de telemetría ocupa entre 60 y 150 bytes, muy por
encima del payload disponible.

\textbf{CayenneLPP (3,55): el competidor más serio en compacidad.} Con 4 bytes por medición y una
implementación trivial resulta imbatible en C-A y C-F. Sin embargo, puntúa 2 en normalización porque
estandariza \emph{tipos de dato} y no \emph{eventos ni estados de salud del dispositivo}: carece de
forma de expresar que un nodo dejó de reportar o que su batería está baja. Y puntúa 2 en
interoperabilidad porque existe exclusivamente dentro del ecosistema LoRaWAN, lo que contradice el
requisito de no cerrarse a una sola tecnología de radio.

\textbf{LwM2M (3,55): técnicamente superior, operativamente desproporcionado.} Es la alternativa más
completa ---modelo de objetos estandarizado, gestión del ciclo de vida del dispositivo, DTLS con
aprovisionamiento inicial de credenciales--- y obtiene 5 en seguridad y telemetría. Su penalización
reside en C-F: la pila completa resulta exigente en memoria de programa para un nodo mínimo, y su
despliegue real se concentra en redes celulares antes que en radio de banda libre. \textbf{Corresponde
reconocer que sería la alternativa preferible si el proyecto contase con mayor presupuesto de hardware y
un único transporte celular.}

\textbf{SIA DC-03 (3,95): la alternativa más cercana.} Es el otro formato publicado por la SIA,
transportado por el mismo DC-09 y con la misma base instalada, de modo que obtiene 5 en normalización
semántica e interoperabilidad. Su vocabulario resulta incluso \textbf{más completo que el de Contact ID
para este proyecto}: incluye el código \texttt{XL} (\emph{Low Received Signal Strength}), que expresa
que la señal de un nodo está débil y para el cual \textbf{Contact ID no tiene equivalente}. Admite
además valores medidos de forma nativa mediante sus bloques de longitud variable.

Se descarta por dos razones. La primera es de eficiencia: sus mensajes son texto ASCII de longitud
variable, del orden de 15 a 25 bytes frente a los 8 de Contact ID en BCD, y componerlos exige
manipulación de cadenas en lugar de concatenación de dígitos, lo que sobre un enlace de radio a batería
se traduce en más tiempo de aire y mayor consumo. La segunda es metodológica: \textbf{la sintaxis exacta
de su bloque de datos para valores numéricos no pudo verificarse}, pues aparece únicamente en el
estándar DC-03, de acceso pagado. Fundar la decisión técnica central del trabajo en un supuesto no
verificado sería incoherente con el criterio de verificación empírica aplicado en el resto del análisis.
Queda registrada como alternativa preferente en caso de que la restricción de compacidad se relaje o de
que se obtenga acceso al estándar (pendiente P8).

\textbf{Contact ID con DC-09 (4,75): gana en cinco de seis criterios y resuelve el sexto mediante un
mecanismo del propio estándar.} El mensaje DC-05 no transporta valores medidos, pero el estándar de
transporte DC-09 define un campo de datos extendidos destinado a información determinada por el
implementador, según se desarrolla en la sección~\ref{sec:extdata}.

\subsubsection{Decisión}

\begin{cajaclave}[title={Protocolo seleccionado}]
\textbf{Se adopta Contact ID (SIA DC-05) sobre transporte DC-09 como protocolo único del perfil.} Los
eventos se expresan en el mensaje DC-05 de 16 dígitos y las mediciones con valor viajan en el
\textbf{campo de datos extendidos que DC-09 define} para información determinada por el implementador
(sección~\ref{sec:extdata}).

La separación entre eventos y telemetría se mantiene como \textbf{separación lógica} ---cada uno tiene
requisitos de entrega distintos--- pero \textbf{no se traduce en dos implementaciones en el
dispositivo}: el nodo emite un solo mensaje.
\end{cajaclave}

Argumento sintético de la decisión:

\begin{enumerate}
\item \textbf{Restricción física:} es el único candidato que combina un mensaje de 8 bytes con un
vocabulario semántico normalizado. Los formatos compactos no normalizan semántica; los que normalizan
semántica no son compactos.
\item \textbf{Alineación con el problema:} su vocabulario ya contiene los eventos de salud de nodos de
radio alimentados a batería ---381, 384, 344, 602 y 603--- que son exactamente las métricas que la
plataforma debe gestionar.
\item \textbf{Interoperabilidad real:} es un estándar ANSI/SIA con base instalada, no un formato propio.
La plataforma puede recibir eventos de paneles de alarma existentes sin modificarlos.
\item \textbf{Disponibilidad de implementación de referencia:} en agosto de 2025 la SIA publicó una
implementación de DC-09 cuyo código es de acceso público, bajo licencia \emph{Polyform Noncommercial},
que \textbf{permite el uso no comercial} ---categoría que comprende este trabajo de titulación--- y
requiere licencia aparte para uso comercial. Constituye por tanto una referencia consultable para
verificar el comportamiento del protocolo sin adquirir la norma, aunque \textbf{no es una dependencia
del proyecto}: la implementación empleada en este trabajo es propia y se verificó de forma
independiente (sección~\ref{sec:alfabeto}).
\item \textbf{Honestidad técnica:} se adopta reconociendo explícitamente el límite del formato DC-05 y
resolviéndolo mediante un mecanismo del propio estándar, no ignorándolo ni inventando una extensión
propietaria.
\end{enumerate}

\subsection{Verificación de la restricción de tamaño}
\label{sec:tamano}

Chile autoriza la banda 915--928\,MHz para uso libre mediante la \textbf{Resolución 737 de Subtel}, y
admite tanto el plan AU915 como el plan US915. La diferencia entre ambos es relevante para el diseño.

\begin{longtable}{@{}L{1.8cm}L{5.2cm}C{1.9cm}C{1.9cm}C{1.9cm}@{}}
\filaencab \encab{Plan} & \encab{Condición} & \encab{DR0} & \encab{DR1} & \encab{DR2} \\ \midrule \endhead
AU915 & Sin restricción de tiempo de permanencia & 51 B & 51 B & 51 B \\
AU915 & Con restricción de tiempo de permanencia & no disp. & no disp. & \textbf{11 B} \\
US915 & Norma FCC, sin restricción & \textbf{11 B} & \textbf{53 B} & 125 B \\
\bottomrule
\caption{Payload de aplicación disponible según plan regional y tasa de datos.}
\label{tab:payload}
\end{longtable}

\begin{cajaaviso}[title={Precisión sobre el límite de 11 bytes}]
La restricción de 11 bytes \textbf{no constituye un techo permanente}: corresponde a la tasa de datos más
lenta, que es la empleada para obtener el alcance máximo. En operación normal el payload disponible es de
51 a 53 bytes. El diseño se dimensiona igualmente contra los 11 bytes, porque un protocolo que solo
funcionara en tasas altas renunciaría al alcance que justifica el empleo de LoRa.
\end{cajaaviso}

Queda pendiente confirmar cuál de los dos planes regionales corresponde emplear en Chile según la
Resolución 737 y si aplica restricción de tiempo de permanencia; la respuesta no altera la viabilidad del
diseño, únicamente el margen disponible.

\begin{longtable}{@{}L{6.4cm}C{3.0cm}L{4.0cm}@{}}
\filaencab \encab{Representación} & \encab{Tamaño} & \encab{Comentario} \\ \midrule \endhead
16 dígitos en ASCII & 16 bytes & No cabe en el peor caso \\
\textbf{16 dígitos en BCD comprimido} & \textbf{8 bytes} & Dos dígitos por byte \\
\bottomrule
\caption{Tamaño del mensaje Contact ID según su representación.}
\label{tab:tamano}
\end{longtable}

\begin{cajaclave}[title={Resultado de la verificación}]
Ocho bytes en un payload de once dejan \textbf{3 bytes de margen}, suficientes para un contador de
secuencia y una marca temporal relativa. El requerimiento RNF-16 queda satisfecho con holgura.
\end{cajaclave}

\begin{longtable}{@{}L{5.2cm}L{4.6cm}L{3.6cm}@{}}
\filaencab \encab{Evento} & \encab{Mensaje} & \encab{BCD (8 bytes)} \\ \midrule \endhead
Nodo 007 dejó de reportar & \texttt{1234 18 1 381 00 007 B} & \texttt{12 34 18 13 81 00 00 7B} \\
Nodo 007 volvió a reportar & \texttt{1234 18 3 381 00 007 9} & \texttt{12 34 18 33 81 00 00 79} \\
Batería baja en nodo 012 & \texttt{1234 18 1 384 00 012 7} & \texttt{12 34 18 13 84 00 01 27} \\
Temperatura alta, zona 003 & \texttt{1234 18 1 158 01 003 7} & \texttt{12 34 18 11 58 01 00 37} \\
Puerta abierta, zona 021 & \texttt{1234 18 1 134 02 021 7} & \texttt{12 34 18 11 34 02 02 17} \\
Reporte periódico & \texttt{1234 18 1 602 00 000 2} & \texttt{12 34 18 16 02 00 00 02} \\
Interferencia RF en gateway & \texttt{1234 18 1 344 00 000 9} & \texttt{12 34 18 13 44 00 00 09} \\
\bottomrule
\caption{Ejemplos de mensaje generados y verificados programáticamente.}
\label{tab:ejemplos}
\end{longtable}

\begin{longtable}{@{}L{6.0cm}C{4.4cm}C{3.0cm}@{}}
\filaencab \encab{Alternativa} & \encab{Tamaño de un evento} & \encab{¿Cabe en 11 B?} \\ \midrule \endhead
\textbf{Contact ID en BCD} & \textbf{8 B} & Sí, con 3 B de margen \\
CayenneLPP (una medición) & 4 B & Sí, pero no expresa el evento \\
CoAP (cabecera y payload mínimo) & 10--20 B & Al límite \\
LwM2M sobre CoAP con CBOR & 20--40 B & No \\
SenML en CBOR & 25--40 B & No \\
Sparkplug B & 40--80 B & No \\
JSON ad hoc & 60--150 B & No \\
\bottomrule
\caption{Comparación del tamaño de mensaje entre alternativas.}
\label{tab:comparaciontamano}
\end{longtable}

\subsection{Limitación de Contact ID y arquitectura de dos planos}
\label{sec:dosplanos}

\begin{cajahallazgo}[title={La limitación, enunciada sin rodeos}]
\textbf{Contact ID transporta eventos discretos codificados, no valores analógicos.} Puede expresar
\enquote{temperatura alta en la zona 003} mediante el código 158, pero \textbf{no puede expresar
\enquote{23,4\,\textdegree C}}. No existe campo alguno para una magnitud medida.
\end{cajahallazgo}

Existen además dos límites de direccionamiento derivados del tamaño de los campos: $15^4 = 50\,625$
cuentas distintas y $15^3 = 3\,375$ zonas por cuenta, reducidas a 999 si se restringe el alfabeto a
dígitos decimales. Para el alcance de este trabajo ambos techos resultan holgados, y se documentan como
límites conocidos del diseño.

\subsubsection{La solución adoptada}

\begin{figure}[htbp]
\centering
\begin{tikzpicture}[font=\footnotesize]

% ---- nodo sensor ----
\node[cajaacento, text width=3.6cm, fill=verdeacento!6] (sensor) at (0,0)
  {\textbf{NODO SENSOR}};

% ---- el mensaje unico ----
\node[caja, below=9mm of sensor, text width=9.6cm, draw=azuloscuro, line width=1.1pt,
      fill=azuloscuro!4, font=\scriptsize] (msg)
  {\normalsize\textbf{\color{azuloscuro}UN SOLO MENSAJE}\\[4pt]
   \texttt{msgdata: 123418138100007B}\qquad \texttt{extdata: [T23.4]}\\[3pt]
   evento Contact ID, 8 B en BCD \ \quad + \quad \ valor en el campo de datos extendidos};

\draw[flechagruesa] (sensor.south) -- (msg.north);
\node[nota, right=1mm of {(0,-0.75)}, anchor=west, font=\scriptsize\itshape]
  {una sola rutina en el firmware};

% ---- plataforma ----
\node[cajaacento, below=9mm of msg, text width=9.6cm, fill=verdeacento!6] (plat)
  {\textbf{PLATAFORMA}\\[2pt]\scriptsize valida el checksum, normaliza y separa por tipo de información};

\draw[flechagruesa] (msg.south) -- (plat.north);

% ---- los dos planos, DESPUES de la plataforma ----
\node[caja, below left=15mm and 3mm of plat.south, anchor=north east,
      text width=5.9cm, align=left, draw=azuloscuro, line width=1pt, font=\scriptsize] (eventos)
  {\textbf{\color{azuloscuro}PLANO DE EVENTOS}\\[3pt]
   vocabulario cerrado de códigos DC-05\\
   máquina de estados: \texttt{Q=1} abre, \texttt{Q=3} cierra\\
   \textbf{entrega confirmada}: un evento perdido\\ es una alarma que no suena};

\node[caja, below right=15mm and 3mm of plat.south, anchor=north west,
      text width=5.9cm, align=left, draw=ambar, line width=1pt, font=\scriptsize] (telemetria)
  {\textbf{\color{ambar}PLANO DE TELEMETRÍA}\\[3pt]
   valor, unidad y marca temporal\\
   serie de tiempo para análisis\\
   \textbf{pérdida tolerable}: una lectura\\ entre mil no altera la tendencia};

\node[nota, below=2mm of plat.south, font=\scriptsize\itshape, fill=white, inner sep=1.5pt] (sep)
  {separación \textbf{lógica}, no de implementación};

\draw[flecha] (plat.south) -- ++(0,-0.95) -| (eventos.north);
\draw[flecha, color=ambar] (plat.south) -- ++(0,-0.95) -| (telemetria.north);

\end{tikzpicture}
\caption{El nodo emite un solo mensaje, con el evento en formato DC-05 y el valor medido en el campo de
datos extendidos que DC-09 reserva para el implementador. La separación en dos planos ---con sus
requisitos de entrega distintos--- ocurre en la plataforma, no en el dispositivo.}
\label{fig:dosplanos}
\end{figure}


\begin{longtable}{@{}L{6.7cm}L{6.7cm}@{}}
\filaencab \encab{Plano de eventos} & \encab{Plano de telemetría} \\ \midrule \endhead
Cruce de umbral (temperatura alta, nivel bajo) & Valor medido periódico \\
Cambio de estado (puerta abierta o cerrada) & Serie de tiempo para análisis \\
Salud del dispositivo (batería, sin supervisión) & Nivel exacto de batería en porcentaje \\
Salud del enlace (interferencia, sin comunicación) & RSSI y SNR numéricos por recepción \\
Confirmación periódica de supervisión & --- \\
Manipulación, avería, anulación & --- \\
\bottomrule
\caption{Criterio de asignación de información a cada plano.}
\label{tab:asignacionplanos}
\end{longtable}

La correlación entre planos se realiza mediante la tripleta de identidad \texttt{(ACCT, GG, CCC)} y la
marca temporal, lo que permite reconstruir el contexto de un evento.

\begin{cajaclave}[title={La separación como decisión de ingeniería}]
Los dos planos presentan requisitos distintos, y separarlos permite tratarlos de manera distinta: el
plano de eventos exige \textbf{entrega confirmada} ---un evento perdido es una alarma que no suena---
mientras que el de telemetría admite \textbf{pérdida tolerable} ---una lectura perdida entre mil no
altera la tendencia. La separación es por tanto una decisión de ingeniería de redes y no una concesión
ante una limitación.
\end{cajaclave}

\subsection{El mecanismo del estándar: campo de datos extendidos de DC-09}
\label{sec:extdata}

La separación en dos planos responde a una necesidad real de diseño ---eventos y mediciones tienen
requisitos de entrega distintos--- pero \textbf{no obliga a que el dispositivo implemente dos protocolos
distintos}. El estándar de transporte DC-09 provee el mecanismo para evitarlo.

\subsubsection{Definición del campo}

DC-09 especifica un campo de datos extendidos con la estructura siguiente:

\begin{center}
\begin{minipage}{0.88\textwidth}
\begin{lstlisting}
[<identificador><datos>]

  identificador : un caracter ASCII mayuscula, de "A" a "Z"
  datos         : caracteres ASCII; el formato lo determina el implementador
  terminacion   : el caracter "]", incluido en el mensaje transmitido
\end{lstlisting}
\end{minipage}
\end{center}

El estándar establece que \textbf{la longitud, el formato y el algoritmo de los datos de este campo los
determina el fabricante}, con una única restricción: no se admite contenido binario.

\subsubsection{Consecuencia para el diseño}

El valor medido no requiere un protocolo aparte: viaja en un campo que el propio estándar reserva para
información definida por el implementador.

\begin{center}
\begin{minipage}{0.88\textwidth}
\begin{lstlisting}
Mensaje DC-09
  token    : *ADM-CID              Contact ID cifrado con AES
  msgdata  : 123418138100007B      el evento, 16 digitos DC-05
  extdata  : [T23.4]               el valor medido
              |+---- dato en ASCII
              +----- identificador asignado por el perfil (T = temperatura)
\end{lstlisting}
\end{minipage}
\end{center}

\begin{longtable}{@{}L{4.4cm}L{9.3cm}@{}}
\filaencab \encab{Propiedad que se conserva} & \encab{Fundamento} \\ \midrule \endhead
No se abandona el estándar & El campo está definido por DC-09; el perfil únicamente determina su contenido \\
La interoperabilidad queda intacta & Una receptora que no interprete el campo lo ignora y procesa el evento igual. Es el comportamiento previsto para datos definidos por el fabricante \\
El dispositivo implementa un solo formato & Una sola rutina de composición de mensaje, con el valor como parámetro opcional \\
\bottomrule
\caption{Propiedades que preserva el empleo del campo de datos extendidos.}
\label{tab:propiedadesextdata}
\end{longtable}

\subsubsection{Realización en cada transporte}

\begin{longtable}{@{}L{3.2cm}L{10.5cm}@{}}
\filaencab \encab{Transporte} & \encab{Cómo viajan el evento y el valor} \\ \midrule \endhead
HTTP y DC-09 sobre IP & Mensaje DC-09 completo, con el evento en \texttt{msgdata} y el valor en \texttt{extdata}. Nativo del estándar \\
LoRaWAN & El nodo no genera tramas DC-09 ---resultan demasiado extensas para el payload disponible---: emite \textbf{un único mensaje} con el evento en 8 bytes y el valor en el mismo paquete. La plataforma compone el mensaje DC-09 cuando corresponde reenviarlo a una receptora externa \\
MQTT & Igual que HTTP, o el formato compacto según la capacidad del cliente \\
\bottomrule
\caption{Realización del mensaje unificado en cada transporte.}
\label{tab:realizacion}
\end{longtable}

\begin{cajaclave}[title={La distinción que importa}]
Los dos planos subsisten como \textbf{separación lógica}, con sus requisitos de entrega diferenciados,
porque esa separación es correcta desde el punto de vista de la ingeniería. Lo que se elimina es la
\textbf{duplicación de implementación}: el dispositivo emite un mensaje y no dos, y quien programa el
firmware implementa una rutina y no dos.

Esta distinción entre separación lógica y unificación de la implementación responde a la objeción
legítima de que dos planos supondrían el doble de trabajo en el nodo. \textbf{No lo suponen.}
\end{cajaclave}

\subsubsection{Identificadores reservados por el perfil}

\begin{longtable}{@{}C{1.6cm}L{7.6cm}L{4.2cm}@{}}
\filaencab \encab{Id.} & \encab{Contenido} & \encab{Ejemplo} \\ \midrule \endhead
\texttt{T} & Temperatura en grados Celsius & \texttt{[T23.4]} \\
\texttt{H} & Humedad relativa en por ciento & \texttt{[H67]} \\
\texttt{N} & Nivel en por ciento o centímetros & \texttt{[N45]} \\
\texttt{B} & Batería en por ciento & \texttt{[B87]} \\
\texttt{V} & Tensión en voltios & \texttt{[V3.28]} \\
\texttt{S} & Calidad de señal: RSSI y SNR & \texttt{[S-103/4.2]} \\
\texttt{Q} & Contador de secuencia del nodo & \texttt{[Q1042]} \\
\bottomrule
\caption{Identificadores del campo de datos extendidos reservados por el perfil CID-IoT. La tabla se
cierra en la Etapa 2.}
\label{tab:identificadores}
\end{longtable}

Queda pendiente determinar si DC-09 impone un límite de longitud a este campo y si existen
identificadores ya asignados por convención que corresponda evitar.


\subsubsection{Verificación del tamaño del mensaje unificado}

El empleo del campo de datos extendidos aumenta el tamaño del mensaje, porque el estándar exige que su
contenido sea texto ASCII y no admite formato binario. Los valores siguientes se obtuvieron ejecutando
la herramienta de verificación del proyecto.

\begin{longtable}{@{}L{5.6cm}C{2.2cm}C{2.2cm}C{2.4cm}@{}}
\filaencab \encab{Mensaje} & \encab{Evento} & \encab{Valor} & \encab{Total} \\ \midrule \endhead
Solo evento (\texttt{381}) & 8 B & --- & \textbf{8 B} \\
Evento con batería (\texttt{[B12]}) & 8 B & 5 B & \textbf{13 B} \\
Evento con temperatura (\texttt{[T23.4]}) & 8 B & 7 B & \textbf{15 B} \\
Evento con RSSI y SNR (\texttt{[S-103/4.2]}) & 8 B & 11 B & \textbf{19 B} \\
\bottomrule
\caption{Tamaño del mensaje unificado según el valor transportado, verificado programáticamente.}
\label{tab:tamanounificado}
\end{longtable}

\begin{cajahallazgo}[title={Consecuencia verificada: el valor debe ser opcional}]
El mensaje unificado ocupa entre 13 y 19 bytes, de modo que \textbf{no cabe en los 11 bytes de la tasa de
datos más lenta}, aunque sí holgadamente en los 51 a 53 bytes de la operación normal.

\textbf{Decisión de diseño derivada:} el valor es \textbf{opcional}. Cuando el enlace opera en la tasa
más lenta ---la de alcance máximo--- el nodo transmite únicamente el evento de 8 bytes, y la medición se
envía cuando el canal lo permite o en un mensaje separado de menor frecuencia.

Esta degradación es deliberada y coherente con la jerarquía de requisitos establecida: \textbf{el evento
es la información que no puede perderse}, porque un evento no entregado es una alarma que no suena,
mientras que una medición ausente solo interrumpe momentáneamente una serie. Cuando el canal se estrecha,
se sacrifica la medición y se preserva el evento.
\end{cajahallazgo}

\subsection{Perfil de aplicación propuesto: CID-IoT}

El aporte original del trabajo no consiste en emplear Contact ID, sino en especificar un \textbf{perfil
de aplicación} que define cómo eventos de sensores heterogéneos se mapean a un vocabulario Contact ID
unificado, y cómo se transportan sobre redes que el protocolo original no contemplaba.

\begin{longtable}{@{}C{2.0cm}L{4.6cm}L{6.7cm}@{}}
\filaencab \encab{Campo} & \encab{Uso original} & \encab{Uso en CID-IoT} \\ \midrule \endhead
\texttt{ACCT} & Número de abonado & \textbf{Sitio u organización} registrada en la plataforma \\
\texttt{GG} & Partición del panel & \textbf{Zona lógica} dentro del sitio \\
\texttt{CCC} & Zona o usuario & \textbf{Dispositivo} dentro de la zona lógica \\
\texttt{MT} & Tipo de mensaje & Se conserva el valor 18, por compatibilidad \\
\texttt{Q} & Calificador & Se conserva: la máquina de estados es funcionalidad deseada \\
\texttt{XYZ} & Código de evento & Se conserva el vocabulario DC-05 \\
\texttt{S} & Checksum & Se conserva el algoritmo módulo 15 \\
\bottomrule
\caption{Asignación del espacio de identidad en el perfil CID-IoT.}
\label{tab:perfil}
\end{longtable}

\begin{cajaclave}[title={Restricción deliberada del perfil}]
\textbf{El perfil no introduce códigos nuevos.} Todo evento se expresa mediante un código DC-05
existente. Esta restricción es intencionada: es lo que preserva la interoperabilidad con las receptoras
instaladas, que constituye el criterio C-C por el cual el protocolo fue seleccionado. Introducir códigos
propios convertiría el estándar en un formato ad hoc y anularía su ventaja.
\end{cajaclave}

\subsubsection{Mapeo de sensores heterogéneos a eventos normalizados}

\begin{longtable}{@{}L{2.6cm}L{5.0cm}L{3.6cm}C{1.2cm}C{0.8cm}@{}}
\filaencab \encab{Transporte} & \encab{Señal del sensor} & \encab{Condición} & \encab{Código} & \encab{Q} \\ \midrule \endhead
LoRaWAN & Sin transmisión en 3 intervalos & Pérdida de supervisión & 381 & 1 \\
LoRaWAN & Reaparece la transmisión & Restauración & 381 & 3 \\
LoRaWAN & Campo de batería bajo umbral & Batería baja de nodo & 384 & 1 \\
LoRaWAN & Transmisión programada recibida & Supervisión activa & 603 & 1 \\
LoRaWAN & Relación señal-ruido baja persistente & Interferencia en receptor & 344 & 1 \\
MQTT & Mensaje de última voluntad recibido & Cliente desconectado & 381 & 1 \\
MQTT & Reconexión del cliente & Restauración & 381 & 3 \\
MQTT / HTTP & Temperatura sobre umbral & Temperatura alta & 158 & 1 \\
MQTT / HTTP & Temperatura vuelve al rango & Restauración & 158 & 3 \\
WiFi / GPIO & Contacto magnético abierto & Apertura de acceso & 134 & 1 \\
WiFi / GPIO & Contacto magnético cerrado & Cierre de acceso & 134 & 3 \\
Cualquiera & Detección de humedad anómala & Fuga de agua & 154 & 1 \\
Cualquiera & Dispositivo en mantenimiento & Zona anulada & 570 & 1 \\
\textbf{Panel de alarma} & \textbf{Ya emite Contact ID nativo} & \textbf{Sin traducción} & --- & --- \\
\bottomrule
\caption{Extracto de la tabla de mapeo del perfil CID-IoT. La tabla completa se desarrolla en la Etapa 2.}
\label{tab:mapeo}
\end{longtable}

La última fila constituye el argumento central del perfil: \textbf{un panel de alarma existente no
requiere traducción, porque ya emplea el protocolo}. La plataforma lo ingiere sin adaptador.

\subsubsection{Vías de ingesta}

\begin{longtable}{@{}C{1.4cm}L{5.0cm}L{7.3cm}@{}}
\filaencab \encab{Vía} & \encab{Transporte} & \encab{Uso} \\ \midrule \endhead
\textbf{A} & Payload LoRaWAN de 8 bytes en BCD, vía integración del servidor de red & Nodos LoRa propios \\
\textbf{B} & DC-09 nativo sobre TCP o UDP con AES & Paneles de alarma e infraestructura instalada \\
\textbf{C} & Interfaz HTTP, mensaje en texto de 16 dígitos & Sensores WiFi, integraciones y pruebas \\
\bottomrule
\caption{Vías de ingesta definidas por el perfil.}
\label{tab:vias}
\end{longtable}

\begin{figure}[htbp]
\centering
\begin{tikzpicture}[font=\scriptsize]

\node[caja, text width=3.9cm, align=left, draw=ambar, line width=1pt] (a) at (0,1.6)
  {\textbf{A}\ \ Nodos LoRa propios\\[2pt]\scriptsize payload de 8 B en BCD, vía el servidor de red};
\node[caja, text width=3.9cm, align=left, draw=verdeacento, line width=1.2pt,
      fill=verdeacento!6] (b) at (0,0)
  {\textbf{B}\ \ Paneles de alarma \textbf{ya instalados}\\[2pt]\scriptsize DC-09 nativo sobre TCP/UDP con AES};
\node[caja, text width=3.9cm, align=left, draw=azuloscuro, line width=1pt] (c) at (0,-1.6)
  {\textbf{C}\ \ Sensores WiFi e integraciones\\[2pt]\scriptsize API HTTP, mensaje de 16 dígitos};

\node[caja, text width=3.3cm, align=center, draw=azuloscuro, line width=1.2pt,
      fill=azuloscuro!5, minimum height=3.2cm] (norm) at (6.6,0)
  {\textbf{NORMALIZACIÓN}\\[4pt]\scriptsize valida el checksum\\ traduce al vocabulario DC-05\\ resuelve la identidad\\ del dispositivo};

\node[caja, text width=2.9cm, align=center, draw=verdeacento, line width=1.1pt,
      fill=verdeacento!6] (uno) at (11.9,0)
  {\textbf{UN SOLO}\\\textbf{ALMACÉN}\\[3pt]\scriptsize formato único\\ normalizado};

\foreach \x in {a,b,c} \draw[flecha] (\x.east) -- (norm.west);
\draw[flechagruesa] (norm.east) -- (uno.west);

\node[notaclave, below=5mm of norm.south, anchor=north, text width=13.2cm, xshift=-0.9cm,
      fill=verdeacento!6, font=\scriptsize]
  {\textbf{La vía B es el diferenciador.} Un panel de alarma existente no requiere traducción alguna,
   porque ya emplea el protocolo: la plataforma lo ingiere sin adaptador. Ningún proveedor regional
   identificado ofrece un punto de ingesta compatible con la infraestructura de monitoreo ya
   desplegada.};

\end{tikzpicture}
\caption{Las tres vías de ingesta convergen a un único formato normalizado. El mismo evento llega por
tres caminos distintos y se almacena una sola vez.}
\label{fig:ingesta}
\end{figure}


\textbf{La vía B constituye el diferenciador}: ningún proveedor regional identificado ofrece un punto de
ingesta compatible con la infraestructura de monitoreo de alarmas ya desplegada.

La vía C responde además a la pregunta sobre sensores WiFi: el mensaje es agnóstico del transporte, de
modo que un microcontrolador con conectividad WiFi construye el mensaje y lo envía por HTTP, sin línea
telefónica, sin módem celular y sin gateway. \textbf{El mismo formato de evento llega por tres caminos
distintos y se almacena una sola vez, normalizado.}

\subsection{Elementos pendientes de especificación}

Los siguientes elementos del perfil se desarrollarán durante la Etapa 2:

\begin{itemize}
\item Gramática formal del mensaje en notación ABNF.
\item Tabla de mapeo completa, con umbral de disparo y de restauración por tipo de sensor.
\item Política de reintento y confirmación, equivalente funcional del \emph{kissoff}.
\item Procedimiento de aprovisionamiento y asignación de identidad.
\item Gestión de claves AES por sitio, con política de rotación alineada a DC-09-2026.
\item Máquina de estados de las condiciones abiertas, con tratamiento de eventos huérfanos.
\item Tratamiento de la desincronización temporal en nodos sin reloj confiable.
\end{itemize}
